\documentclass[11pt]{article}
\usepackage[T1]{fontenc}
\usepackage[utf8]{inputenc}
\usepackage{lmodern}
\usepackage[english]{babel}
\usepackage{amsmath,amssymb,mathtools,array,booktabs,pdflscape,adjustbox,diagbox}
\usepackage[a4paper,margin=1.45cm]{geometry}
\usepackage[protrusion=true,expansion=false]{microtype}
\setlength{\parindent}{1.25em}
\setlength{\parskip}{0pt}
\allowdisplaybreaks
\setlength{\emergencystretch}{12em}
\sloppy
\hbadness=10000
\vbadness=10000
\hfuzz=2pt
\vfuzz=10pt
\raggedbottom
\makeatletter
\renewcommand\footnotesize{\@setfontsize\footnotesize{7.5}{8.5}\selectfont}
\makeatother
\providecommand{\frakA}{\mathfrak A}
\providecommand{\frakB}{\mathfrak B}
\providecommand{\frakC}{\mathfrak C}
\providecommand{\frakM}{\mathfrak M}
\providecommand{\frakN}{\mathfrak N}
\providecommand{\frako}{\mathfrak o}
\providecommand{\fraka}{\mathfrak a}
\providecommand{\frakc}{\mathfrak c}
\providecommand{\frakp}{\mathfrak p}
\providecommand{\frakO}{\mathfrak O}
\providecommand{\frakP}{\mathfrak P}
\providecommand{\frakD}{\mathfrak D}
\providecommand{\frakG}{\mathfrak G}
\providecommand{\frakK}{\mathfrak K}
\providecommand{\Om}{\Omega}
\providecommand{\N}{\operatorname{N}}
\providecommand{\Sp}{\operatorname{Sp}}
\providecommand{\rank}{\operatorname{rank}}
\providecommand{\Trdeg}{\operatorname{Trdeg}}
\providecommand{\charac}{\operatorname{char}}
\begin{document}
\section*{36. Ideal Differentiation and the Different}

\begin{center}
\emph{J. Ber. d. DMV 39 (1930), p. 17}
\end{center}

E. Noether, Göttingen: Ideal differentiation and the different.

It is shown how the different of an algebraic number field can be regarded as the differential quotient of a defining ideal. That this definition agrees with the usual one follows from structural theorems which are interesting in themselves and which, in an analogous sense, constitute a generalization of the Lagrange interpolation formula.

A detailed account is to appear in the \emph{Mathematische Annalen}.

\section*{37. Normal Basis in Fields without Higher Ramification}

\begin{center}
\emph{Journal f. d. reine u. angew. Math. 167 (1932), pp. 147--152}
\end{center}

By a normal basis of a Galois number field \(K/k\) one means a basis of the principal order \(\frakO\) of \(K\), relative to the principal order \(\frako\) of \(k\), consisting of conjugate elements. A necessary condition for the existence of such a basis is, as is known, that no higher ramification fields occur in \(K/k\). This remains true even if one restricts to a place, i.e. passes to the \(p\)-adic extension \(k_p\) of \(k\) and the corresponding \(K_p\) of \(K\). I prove below the converse:

For every place \(\frakp\) not dividing the degree \(n\) of \(K/k\), a normal basis exists. In particular: if \(K/k\) has no higher ramification, then at every ramified place a normal basis exists; the discriminant there is the square of a group determinant.

To this last assertion one must add that, contrary to a conjecture of E. Artin, the passage to his conductors still requires further considerations: decomposing the group determinant according to irreducible characters yields generalized root numbers; a suitable grouping yields a factorization in the ground field enlarged by the characters. In the simplest cases I can identify these new factors with Artin's conductors by means of the ideal-theoretic decomposition of the root numbers. In the general case, where no normal basis exists, a splitting into generalized root numbers is still always possible by means of a composition series by Galois modules; from this one again obtains a corresponding decomposition in the enlarged ground field.

The proof of the existence of the normal basis rests on the fact that \(\frakO/\frako\), viewed as a Galois module, becomes operator-isomorphic to a one-sided ideal of the integral group ring extended by \(\frako\). At all ordinary ramification places this group ring is a maximal order; ideals in maximal orders of \(p\)-adic semisimple systems are principal ideals. Returning from this to the Galois module gives precisely the existence of the normal basis. At the higher ramification places, the integral group ring passes to a nonmaximal order, and the ideal corresponding to \(\frakO/\frako\) becomes a nonprincipal ideal.

All these considerations remain valid for function fields of one variable, provided that the characteristic of the constant field does not divide the degree of \(K/k\).

\subsection*{§1. The \(p\)-adically Extended Integral Group Ring}

Let \(\frakO\) and \(\frako_p\) denote the principal orders of an algebraic number field \(k\) and of its \(p\)-adic extension \(k_p\), where \(p\) is a prime ideal of \(k\). Let \((G)\) denote the rational integral group ring and \([G]\) the integral group ring of a group \(G\) with \(n\) elements. By \((G)_k\), respectively \((G)_{k_p}\), we mean the group ring with coefficients in \(k\), respectively \(k_p\); analogously \([G]_{\frako}\) and \([G]_{\frako_p}\) are the corresponding extensions of the integral group ring to coefficients in \(\frako\), respectively \(\frako_p\).

\paragraph{Theorem 1.}
If \(p\) does not divide the number \(n\) of elements of \(G\), then \([G]_{\frako_p}\) is a maximal order of \((G)_{k_p}\). If \(p\) divides \(n\), then \([G]_{\frako_p}\) is not maximal.

The discriminant of the integral group ring, and hence of \([G]_{\frako}\) and \([G]_{\frako_p}\), is a power of \(n\). Thus if \(p\) does not divide \(n\), the discriminant of \([G]_{\frako_p}\) is a unit. This means that, while \(\frako_p\) is fixed, \([G]_{\frako_p}\) cannot be enlarged to a larger order; such an enlargement would split off a nonunit square factor from the discriminant. This proves the first part.

If, on the other hand, \(p\) divides \(n\), the direct sum corresponding to the identity representation,
\[
        E^0 [G]_{\frako_p}+(1-E^0)[G]_{\frako_p},
        \qquad
        E^0=\frac1n\sum_{S\in G}S,
\]
is a proper extension of \([G]_{\frako_p}\).

\paragraph{Theorem 2.}
If \(p\) does not divide \(n\), every integral or fractional ideal with respect to \([G]_{\frako_p}\) is principal. Indeed, by Theorem 1, \([G]_{\frako_p}\) is a maximal order in a \(p\)-adic semisimple system, and the ideals of such orders are principal.

\subsection*{§2. Galois Modules, Operator Isomorphism, Normal Basis}

Let \(K/k\) be Galois with group \(G\). For \(z\in K\), let \(z^S\) denote the element obtained from \(z\) by the substitution \(S\in G\). The substitutions in \(G\) induce operator automorphisms of \(K/k\), regarded as a \(k\)-module. Thus \(K/k\) can be made into a module with respect to \((G)_k\) by setting
\[
        z\Bigl(\sum_i c_iS_i\Bigr)=\sum_i c_i z^{S_i},
        \qquad z\in K,\; c_i\in k,\; S_i\in G.
\]

\paragraph{Definition.}
Every \((G)_k\)-module contained in \(K/k\) is called a rational Galois module. Similarly, the \([G]_{\frako}\)-modules contained in \(K/k\) are called integral Galois modules; in particular \(\frakO/\frako\) is an integral Galois module.

\paragraph{Theorem 3.}
As a rational Galois module, \(K/k\) is operator-isomorphic to \((G)_k\). This follows from the fact that \(K/k\) always has a field normal basis, i.e. a \(k\)-basis consisting of conjugates \(z^S\). The correspondence
\[
        \sum_i c_iS_i\longmapsto \sum_i c_i z^{S_i}
\]
is an operator homomorphism, and since the ranks are equal it is an isomorphism.

\paragraph{Theorem 4.}
As an integral Galois module, \(\frakO/\frako\) is operator-isomorphic to a \([G]_{\frako}\)-module in \((G)_k\) of maximum rank \(n\). Likewise, \(\frakO_p/\frako_p\) is operator-isomorphic to a \([G]_{\frako_p}\)-module of rank \(n\) in \((G)_{k_p}\).

The isomorphism of Theorem 3 assigns to the \([G]_{\frako}\)-module \(\frakO/\frako\) a \([G]_{\frako}\)-module of rank \(n\) in \((G)_k\). Locally one extends coefficients from \(k\) to \(k_p\); the resulting system is semisimple, in general a sum of isomorphic fields.

\paragraph{Theorem 5.}
At every place \(p\) which does not divide the degree \(n\) of \(K/k\), the module \(\frakO_p/\frako_p\) has a normal basis.

By Theorem 1, \([G]_{\frako_p}\) is a maximal order there; the \([G]_{\frako_p}\)-modules of rank \(n\) in \((G)_{k_p}\) are integral or fractional ideals, hence principal by Theorem 2. If the ideal corresponding to \(\frakO_p/\frako_p\) is \(W[G]_{\frako_p}\), then it has the \(\frako_p\)-basis \(WS_i\). The operator isomorphism says that \(w^{S_i}\) is an \(\frako_p\)-basis of \(\frakO_p/\frako_p\), where \(w\) is the element corresponding to \(W\). Thus the conjugate elements \(w^{S_i}\) form a normal basis.

If \(p\) divides \(n\), the corresponding module cannot be principal, since in this case no normal basis can exist.

\paragraph{Additional remark to Theorem 3.}
From the operator isomorphism between \((G)_k\) and \(K/k\) follow Speiser's results on Klein's problem of forms: every representation of the Galois group \(G\) of \(K/k\) which is irreducible over \(k\) is generated by an irreducible rational Galois module from \(K/k\); every absolutely irreducible representation is generated by a Galois module from \(K_Z/Z\), where \(Z\) is a splitting field of the corresponding component of the group ring. If \(v_1,\ldots,v_m\) is a basis of such a Galois module and \(S\mapsto (s_{ij})\) is the corresponding representation, then
\[
        v_i^S=\sum_j s_{ij}v_j.
\]
This is precisely Klein and Speiser's formulation.

\subsection*{§3. The Discriminant as Group Determinant in Fields without Higher Ramification}

For decomposing the discriminant it is enough to consider its decomposition at the individual ramified places. In fields without higher ramification, only places with normal basis are involved.

\paragraph{Theorem 6.}
For Galois fields \(K/k\) without higher ramification, the discriminant at every ramified place is the square of the group determinant of the Galois group after replacing the indeterminates by a normal basis. Corresponding to the irreducible representations, the group determinant decomposes into generalized root numbers; the discriminant decomposes into ideals of the ground field enlarged by the characters, with conjugate complex characters corresponding to the same ideals.

With \(w^S\) as a normal basis, the discriminant \(\Delta\) is the square of
\[
        D=\det\bigl(w^{ST}\bigr)_{S,T\in G}.
\]
This is the group determinant with the \(w^S\) in place of the indeterminates \(x_S\). If
\[
        D=D_1\cdots D_h
\]
is the factorization according to the absolutely irreducible representations, then for a representation \(S\mapsto \rho_i(S)\)
\[
        D_i=\det\Bigl(\sum_{S\in G} w^S\rho_i(S)\Bigr).
\]
Thus the \(D_i\) are generalized root numbers. For cyclic groups they are simply the Lagrange root numbers
\[
        \sum_{S\in G} w^S\chi_i(S).
\]

To pass from the factorization of the group determinant to that of the discriminant, one groups each representation with its adjoint. If \(D_i\) and \(D_i'\) are the corresponding factors, put
\[
        \Delta_i=D_iD_i'.
\]
Then \(\Delta_i^T=\Delta_i\) for all \(T\in G\); hence \(\Delta_i\) lies in the ground field enlarged by the real subfield of the \(i\)-th character. The discriminant factorization
\[
        \Delta=\Delta_1\cdots\Delta_h
\]
is therefore a factorization in the field obtained by adjoining the real subfields of the characters, with conjugate complex characters corresponding to the same factors.

If one can show that the ideals generated by the \(\Delta_i\) actually already lie in the base field \(k\), and are equal for conjugate characters, then they agree with Artin's conductors, provided that to composite characters one assigns the corresponding products of powers of the \(\Delta_i\).

\paragraph{Theorem 7.}
Let \(k\) be the field of rational numbers, and let \(K/k\) be cyclic of prime degree \(l\), prime to the discriminant, hence without higher ramification. Then the ideals generated by the \(\Delta_\lambda\) for the nontrivial representations are all equal and agree with the conductor of \(K/k\).

This follows from the known decomposition of root numbers. At a ramified place \(p\) of \(K/k\), for \(k_\varepsilon\), the field of \(l\)-th roots of unity, and \(K_\varepsilon\), one has
\[
        (p)=\mathfrak p_1\cdots\mathfrak p_{l-1},\qquad
        \mathfrak p_i=\mathfrak P_i^{\,l},
\]
with \(\mathfrak p_i\) in \(k_\varepsilon\) and \(\mathfrak P_i\) in \(K_\varepsilon\). Further, for the root numbers,
\[
        (\Omega_\lambda)=\mathfrak P_1^{r_1}\cdots
        \mathfrak P_{l-1}^{r_{l-1}},
        \qquad
        (\Omega_{\bar\lambda})=
        \mathfrak P_1^{l-r_1}\cdots \mathfrak P_{l-1}^{l-r_{l-1}},
\]
where \(r_1,\ldots,r_{l-1}\) is a permutation of \(1,\ldots,l-1\). The root numbers \(\Omega_\lambda\) are here identical with the factors \(D_\lambda\) of the group determinant; hence
\[
        (\Delta_\lambda)=(D_\lambda D_{\bar\lambda})
        =(\Omega_\lambda\Omega_{\bar\lambda})
        =\mathfrak P_1^l\cdots\mathfrak P_{l-1}^l
        =\mathfrak p_1\cdots\mathfrak p_{l-1}
        =(p),
\]
and therefore there is agreement with the conductor of \(K/k\).

\begin{center}
Received 24 August 1931.
\end{center}

\section*{38. Jointly with R. Brauer and H. Hasse: Proof of a Main Theorem in the Theory of Algebras}

\begin{center}
\emph{Journal f. d. reine u. angew. Math. 167 (1932), pp. 399--404}
\end{center}

At last our joint efforts have succeeded in proving the following theorem, which is of fundamental importance for the structure theory of algebras over algebraic number fields and beyond.

\paragraph{Main Theorem.}
Every normal division algebra over an algebraic number field is cyclic, or, equivalently, of Dickson type.

It is a special pleasure to present this result, essentially an achievement of the \(p\)-adic method, to Kurt Hensel, the founder of that method, on his seventieth birthday. Our proof consists of three reductions, each contributed by one of us.

H. Hasse gave the first reduction on the basis of his recently developed theory of cyclic algebras over algebraic number fields.

\paragraph{Reduction 1.}
The main theorem is proved once the following is shown:
\[
\text{I. Every everywhere split algebra over }\Om\text{ is } \sim\Om .
\]
Here and below, an ``algebra \(A\) over \(\Om\)'' means a normal simple algebra over the algebraic number field \(\Om\), i.e. a simple hypercomplex system with center \(\Om\). The notation \(A\sim\Om\) means that \(A\) is a full matrix algebra over \(\Om\), equivalently that it belongs to the special similarity class determined by \(\Om\) itself as division algebra. An algebra is called everywhere split if for every prime place \(\mathfrak p\) of \(\Om\) its \(\mathfrak p\)-adic extension \(A_{\mathfrak p}\) is split:
\[
        A_{\mathfrak p}\sim\Om_{\mathfrak p}.
\]

Let \(D\) be a division algebra over \(\Om\). Choose a cyclic field \(Z/\Om\) such that, for every prime place \(\mathfrak p\), the \(\mathfrak p\)-degree of \(Z\) is a multiple of the \(\mathfrak p\)-index of \(D\). Then for every prime divisor \(\mathfrak P\) of \(\mathfrak p\) in \(Z\), the local field \(Z_{\mathfrak P}\) splits \(D_{\mathfrak p}\). Hence the algebra \(D_Z\) obtained by extending scalars from \(\Om\) to \(Z\) is everywhere split. If I is known, \(D_Z\sim Z\). Thus \(D\) has a cyclic splitting field \(Z\), hence is cyclically representable; moreover it has such a cyclic splitting field whose degree agrees with the degree of \(D\). Therefore \(D\) is cyclic.

The second reduction was prompted decisively by R. Brauer, who informed H. Hasse by letter that the related question of the exact value of the exponent of a division algebra can be reduced, using Sylow's theorem, to the case of a solvable splitting field. The same idea gives the next reduction of the cyclicity question.

\paragraph{Reduction 2.}
Statement I is proved once the following is shown:
\[
\text{II. Every everywhere split algebra over }\Om\text{ with a solvable splitting field is } \sim\Om .
\]
Let \(A\) be everywhere split over \(\Om\). Let \(K\) be a Galois splitting field for \(A\), let \(p\) be a prime number, and let \(\Sigma\) be the Sylow field of \(K/\Om\) corresponding to \(p\), i.e. the fixed field of a \(p\)-Sylow subgroup of the Galois group. Then \(A_\Sigma\) is an everywhere split algebra over \(\Sigma\), and has the solvable splitting field \(K\). By II, \(A_\Sigma\sim\Sigma\). Thus \(\Sigma\) is a splitting field of \(A\), and its degree over \(\Om\) is prime to \(p\). Hence the index of \(A\), being a divisor of that degree, is prime to \(p\). Since this holds for every prime \(p\), the index is \(1\), and \(A\sim\Om\).

The third reduction, which together with the first two leads to the final proof, was extracted by E. Noether from a communication of Hasse; independently, Brauer had also considered it.

\paragraph{Reduction 3.}
Statement II is proved once the following is shown:
\[
\text{III. Every everywhere split algebra over }\Om\text{ with a cyclic splitting field of prime degree is } \sim\Om .
\]
Let \(A\) be everywhere split and have a solvable splitting field \(K/\Om\). Choose a tower
\[
        K=\Lambda_0>\Lambda_1>\cdots>\Lambda_r=\Om
\]
such that each \(\Lambda_i/\Lambda_{i+1}\) is cyclic of prime degree. Then \(A_{\Lambda_1}\) is everywhere split over \(\Lambda_1\), and has the cyclic splitting field \(K=\Lambda_0\) of prime degree. By III, \(A_{\Lambda_1}\sim\Lambda_1\), so \(\Lambda_1\) already splits \(A\). Repeating the same argument down the tower gives \(A\sim\Om\). Thus II follows from III.

But III is known from Hasse's results. Hence the main theorem follows by reversing the reductions 3, 2, and 1. Noether further observes that III, in the more general form for cyclic splitting fields of arbitrary degree, is equivalent to Hasse's norm theorem. For Reduction 3 one needs only the prime-degree case, the original Hilbert--Furtwängler norm theorem. Thus Reduction 3 gives a new simple proof of Hasse's norm theorem. In Hasse's notation: if \(Z/\Om\) is cyclic and \(\alpha\in\Om\) satisfies
\[
        \left(\frac{\alpha,Z}{\mathfrak p}\right)=1
\]
for every prime place \(\mathfrak p\), then every cyclic algebra
\[
        A=(\alpha,Z)
\]
is everywhere split. Reduction 3 gives \(A\sim\Om\), and hence \(\alpha\) is the norm of an element of \(Z\).

\subsection*{Consequences}

\paragraph{Theorem 1.}
The exponent of a normal simple algebra over an algebraic number field is equal to its index.

\paragraph{Theorem 2.}
The fundamental ideal of a proper normal division algebra over \(\Om\) is divisible by at least one prime place of \(\Om\), and in fact by at least two.

The fundamental ideal may be defined as the reduced norm of the reduced different. An infinite prime place is included exactly when the algebra reduces there to the quaternion algebra, rather than to the real or complex number field. By Hasse's results, a prime place occurs in the fundamental ideal exactly when the \(\mathfrak p\)-index is not \(1\). If all \(\mathfrak p\)-indices were \(1\), the algebra would be everywhere split; then by I it would not be a proper division algebra. Nor can a single \(\mathfrak p\)-index alone be different from \(1\), because the norm residue symbols, whose orders are the local indices, satisfy the product theorem.

\paragraph{Theorem 3.}
For an algebra \(A\) over \(\Om\), an algebraic number field \(K/\Om\) is a splitting field if and only if, for every prime divisor \(\mathfrak P_i\) in \(K\) of every prime place \(\mathfrak p\) of \(\Om\), the \(\mathfrak P_i\)-degree \(n_{\mathfrak P_i}\) of \(K\) is a multiple of the \(\mathfrak p\)-index \(m_{\mathfrak p}\) of \(A\).

If
\[
        \mathfrak p=\mathfrak P_1^{e_{\mathfrak P_1}}\cdots
        \mathfrak P_g^{e_{\mathfrak P_g}},
        \qquad
        \N(\mathfrak P_i)=\mathfrak p^{\,f_{\mathfrak P_i}},
\]
then
\[
        n_{\mathfrak P_i}=e_{\mathfrak P_i}f_{\mathfrak P_i}.
\]
The \(\mathfrak p\)-index of \(A\) is the index \(m_{\mathfrak p}\) of \(A_{\mathfrak p}\).

The assertion follows because \(K\) splits \(A\) precisely when \(A_{K_{\mathfrak P_i}}\sim K_{\mathfrak P_i}\) for all \(\mathfrak P_i\). At a finite prime place let
\[
        A_{\mathfrak p}\sim(\alpha,W_{\mathfrak p})
\]
be the distinguished cyclic local representation, where \(W_{\mathfrak p}\) is the unramified field of degree \(m_{\mathfrak p}\) over \(\Om_{\mathfrak p}\). Let \(B_i\) be the intersection and \(C_i\) the composite of \(W_{\mathfrak p}\) and \(K_{\mathfrak P_i}\). With
\[
        d_i=(m_{\mathfrak p},f_{\mathfrak P_i}),
\]
the extension \(B_i/\Om_{\mathfrak p}\) has degree \(d_i\), and \(K_{\mathfrak P_i}/B_i\) is unramified of degree \(m_{\mathfrak p}/d_i\). The local algebra splits exactly when \(\alpha\) is a norm from \(K_{\mathfrak P_i}\) over \(B_i\), which is equivalent to the valuation of \(\alpha\) in \(\mathfrak P_i\) being a multiple of \(m_{\mathfrak p}/d_i\). Since
\[
        \left(\frac{m_{\mathfrak p}}{d_i},\frac{f_{\mathfrak P_i}}{d_i}\right)=1,
\]
this is equivalent to \(n_{\mathfrak P_i}\) being a multiple of \(m_{\mathfrak p}\). For infinite prime places, unless \(m_{\mathfrak p}=1\), one has \(m_{\mathfrak p}=2\), and the condition says precisely that the corresponding local field is complex.

\paragraph{Theorem 4.}
\emph{Decomposition theorem.} Let \(K/\Om\) be Galois, and let \(\mathfrak p\) be a prime ideal of \(\Om\) not dividing the relative discriminant of \(K\). The relative degree \(f\) of the prime divisors of \(\mathfrak p\) in \(K\) is the least exponent for which
\[
        A^f\sim\Om
\]
for every algebra \(A\) in the subgroup of Brauer's group assigned to \(K\).

Indeed, \(f\) is the local degree of \(K\) at \(\mathfrak p\), while the \(\mathfrak p\)-index of \(A\) is the index, equivalently the exponent, of \(A_{\mathfrak p}\); Theorem 3 supplies the local criterion. The existence of algebras with exact \(\mathfrak p\)-index \(f\) follows from Frobenius' density theorem and the construction of a cyclic algebra \((\alpha,Z)\) whose norm residue symbols at two suitably chosen prime places have reciprocal values of order \(f\) and are trivial elsewhere.

\paragraph{Theorem 5.}
\emph{Uniqueness and ordering theorem.} \(K\subseteq K'\) if and only if \(\mathfrak R_K\subseteq\mathfrak R_{K'}\). In particular, the correspondence between Galois fields \(K\) and their algebra groups is one-to-one and order-preserving.

The implication \(K\subseteq K'\Rightarrow\mathfrak R_K\subseteq\mathfrak R_{K'}\) is immediate. Conversely, the decomposition theorem gives the required divisibilities of relative degrees for all primes outside the relative discriminant; the standard analytic argument then yields \(K\subseteq K'\).

\paragraph{Theorem 6.}
All absolutely irreducible representations of a finite group \(G\) can be realized over cyclotomic fields; in any case they can be realized over the field of \(n^h\)-th roots of unity for \(n=|G|\) and sufficiently large \(h\).

Passing to the rational integral group ring of \(G\), the absolutely irreducible representations of \(G\) are the absolutely irreducible representations of the simple components \(G_i\) of that semisimple algebra. Their centers are the fields of the corresponding characters, hence cyclotomic fields. By Theorem 3 it suffices, at the finitely many prime divisors of the fundamental ideal, to make the local degrees divisible by the corresponding indices. The field of \(n^h\)-th roots of unity does this for sufficiently large \(h\).

\begin{center}
Received 11 November 1931.
\end{center}

\section*{39. Hypercomplex Systems in Their Relations to Commutative Algebra and Number Theory}

\begin{center}
\emph{Proceedings of the International Congress of Mathematicians, Zürich 1 (1932), pp. 189--194}
\end{center}

\paragraph{1.}
The theory of hypercomplex systems, or algebras, has advanced strongly in recent years; but only very recently has the significance of this theory for commutative questions become clear. I want to report on this significance of the noncommutative for the commutative, and I shall do so by following two classical questions going back to Gauss: the principal genus theorem and the closely related norm theorem. Over time these questions have repeatedly changed their formulation. In Gauss they appear as the conclusion of the theory of quadratic forms; then they play an essential role in the characterization of relatively cyclic and abelian number fields by class field theory; finally they can be stated as theorems on automorphisms and on the splitting of algebras. This last formulation at the same time transfers the theorems to arbitrary relatively Galois number fields.

The principle of applying the noncommutative to the commutative is this: by means of the theory of algebras one seeks invariant and simple formulations for known facts about quadratic forms or cyclic fields, formulations depending only on structural properties of algebras. Once such invariant formulations are proved, the transfer to arbitrary Galois fields is obtained automatically.

\paragraph{2.}
First, an overview of the methods and further results. The main difficulty lies in finding the right formulation for general Galois fields; without the hypercomplex method there was no starting point for this. In the examples mentioned above, the underlying passage to the noncommutative is obtained by considering field and group simultaneously through the \emph{crossed product} and its multiplication constants, the \emph{factor systems}. Such crossed products were first considered by Dickson, whereas the theory of factor systems was developed by Speiser, Schur, and R. Brauer from a quite different starting point, namely the question of absolutely irreducible representations. Only by combining the two theories did one obtain a sufficiently simple and far-reaching structure.

At the same time this gives new and transparent proofs of known facts: Hasse's hypercomplex proof of the reciprocity law for cyclic fields by means of an invariant formulation of his norm residue symbol; and Chevalley's hypercomplex foundation of local class field theory on the same basis, for which new algebraic theorems on factor systems were needed. One must add the limitation that the crossed-product method by itself apparently does not yield the full theory of Galois number fields: new results of Artin, proceeding from Hasse's proof in the sense of this principle, give only numerical equalities in place of full isomorphism theorems.

Methods giving a full isomorphism, namely operator isomorphism, already exist in algebra. They continue ideas of A. Speiser: the Galois field is viewed as a \emph{Galois module}, i.e. as a module over the base field which admits the substitutions of the Galois group as operators. There is an operator isomorphism between the field and the group ring in the sense that elements correspond one-to-one, linear forms over the base field correspond, and the substitutions of the Galois group in the field correspond to multiplication in the group ring. A theorem formulated by me was proved by M. Deuring, who based on it a proof of Galois theory. Further structural theorems of Deuring run parallel to formal facts of Artin \(L\)-series and give a structural approach to Artin conductors. These Artin \(L\)-series and conductors, formed with general group characters, together with Speiser's work, form the first connection between number theory and representation theory.

\paragraph{3.}
I now examine the norm theorem and the principal genus theorem in detail. First, the definition of the crossed product. Let \(K/k\) be a Galois field of degree \(n\), with group \(G\). The crossed product means a simultaneous embedding of \(K\) and \(G\) in an algebra \(A\) in such a way that the automorphisms of \(K\) become inner. Let \(u_S\) denote symbols corresponding to the \(n\) group elements. First set up \(A\) as a module of linear forms of rank \(n\) over \(K\):
\[
        A=u_{S_1}K+\cdots+u_{S_n}K ,
\]
i.e. \(A\) consists of all linear forms \(u_{S_1}a_1+\cdots+u_{S_n}a_n\), \(a_i\in K\). The inner automorphism condition gives
\[
        u_S^{-1}zu_S=z^S,\qquad\text{or}\qquad
        zu_S=u_Sz^S\quad (z\in K),
\]
and the multiplication is determined by
\[
        u_Su_T=u_{ST}a_{S,T},\qquad a_{S,T}\in K^* .
\]
Associativity gives the factor-system condition
\[
        a_{S,T}a_{ST,R}=a_{S,TR}\,a_{T,R}^{\,S}.
\]
The resulting \(A\) is called the crossed product of \(K\) with \(G\) with factor system \(a_{S,T}\). One proves that \(A\) is a normal simple algebra over \(k\), that \(K\) is a maximal commutative subfield, hence a splitting field, and conversely that for a given division algebra \(D\) there are always matrix rings \(D_r\) that can be obtained in this way as crossed products.

Replacing \(u_S\) by \(v_S=u_Sc_S\), \(c_S\in K^*\), gives associated factor systems
\[
        a'_{S,T}=a_{S,T}\frac{c_S^{\,T}c_T}{c_{ST}}.
\]
Associated factor systems are grouped into a class \((a)\); similarly, all algebras similar to \(A\) are grouped into a class \(\mathfrak A\). The classes \(\mathfrak A\) and \((a)\) correspond bijectively. The classes with fixed splitting field \(K\) form, under direct product, an abelian group isomorphic to the componentwise product of classes of factor systems. The identity element is the split algebra class, respectively the system of all transformation quantities
\[
        \frac{c_S^{\,T}c_T}{c_{ST}}.
\]

\paragraph{4.}
For cyclic splitting fields this specializes to the norm concept. Let \(Z/k\) be cyclic, with \(S\) a generator of its group. One may let the powers of \(u\) correspond to the powers of \(S\):
\[
        A=Z+uZ+\cdots+u^{n-1}Z,
\]
\[
        zu=uz^S,\qquad u^n=a,\qquad a\in k^*,
\]
and, if \(v=uc\),
\[
        a'=a\,N(c).
\]
Thus each factor system is represented by a single element \(a\in k^*\), and one writes
\[
        A=(a,Z).
\]
The identity class of factor systems is given by the norms from \(Z^*\), so the group of algebra classes is isomorphic to
\[
        k^*/N(Z^*).
\]
A cyclic algebra \((a,Z)\) therefore splits if and only if \(a\) is the norm of an element of \(Z\).

This connection between norm and splitting gives the formulation of the norm theorem as a theorem on split algebras: if an algebra splits at every place, then it splits absolutely. Here a place is defined, as in number theory, by replacing the base field \(k\) with its \(p\)-adic extension, and at infinite places by extension to the real numbers. For cyclic algebras this theorem says: if \(a\) is a local norm at every finite and infinite place, then \(a\) is the norm of a number of \(Z\); equivalently, without passing to \(p\)-adic fields, if \(a\) is a norm residue modulo every prime ideal and satisfies the known sign conditions, then \(a\) is a global norm. This is precisely the norm theorem of class field theory. The general theorem on split algebras follows from this cyclic special case by purely algebraic-arithmetic arguments. One first important consequence is Hasse's theorem: every simple normal algebra over an algebraic number field is cyclic.

\paragraph{5.}
A second consequence is the principal genus theorem. Its invariant formulation rests on the fact that the defining relations of the crossed product are purely multiplicative, and therefore still make sense when \(K^*\) is replaced by an abelian group \(\mathfrak G\) whose automorphism group contains a subgroup isomorphic to \(G\). The module of linear forms is then replaced by the extension of \(G\) by \(\mathfrak G\) in the sense of group theory. If \(\mathfrak G\) is taken to be the group of all ideals of \(K\), the factor systems become systems of ideals. A division into classes in \(\mathfrak G\) induces a division into classes for the factor systems.

The identity class of factor systems is defined to consist of those elements \(a_{S,T}\) which generate split algebras at all finite and infinite ramified places of \(K\). Let the corresponding extension of \(G\) be denoted by \(G^*\). Then the invariant formulation of the principal genus theorem is:

If, under this class division, the substitutions
\[
        v_S=u_S(c_S)
\]
produce an automorphism of \(G^*\), that is, if the transformation quantities
\[
        \frac{(c_S)(c_T)^S}{(c_{ST})}
\]
belong to the identity ideal class of factor systems, then there is an ideal class \((\mathfrak b)\) such that
\[
        (c_S)=\frac{(\mathfrak b)}{(\mathfrak b)^S}
\]
for all \(S\in G\). The hypothesis expresses the automorphism property; the fact that this automorphism is inner is expressed by
\[
        v_S=(\mathfrak b)^{-1}u_S(\mathfrak b)
        =u_S(\mathfrak b)^{\,1-S}=u_S(c_S).
\]
In the cyclic special case one obtains: if \(N([c])\) lies in the identity ideal class of factor systems, then \((c)\) is a symbolic \((1-S)\)-st power,
\[
        (c)=(\mathfrak b)^{\,1-S}.
\]

\paragraph{6.}
The theorem has been stated for full ideal classes; its content is unchanged if, as usual, one restricts to ideals prime to the ramified places. With this restriction, the principal genus defined here for quadratic fields becomes Gauss's principal genus: ideal classes correspond to quadratic forms, and the norms of classes correspond to the numbers represented by the forms. The identity class means that these represented numbers are quadratic residues at the ramified places; the associated forms therefore have the total character of the principal form and constitute Gauss's principal genus. It is known that the symbolic \((1-S)\)-st power becomes duplication.

For cyclic fields the statement becomes: the principal genus consists of all ideal classes whose norms are norm residues at the finite and infinite ramified places. This is equivalent to being a norm residue modulo the conductor, and the usual theorem is obtained as a specialization.

Even for arbitrary Galois fields one can introduce a conductor composed only of the ramified places, in such a way that after normalizing the factor systems, at each such place the ray contains only elements of the identity class. This raises the question of the connection with Artin conductors, and hence with the theory of Galois modules, the second hypercomplex method. The future must show how far these two methods reach.
\end{document}
